InsuLink is a mobile application for people living with diabetes that goes beyond displaying glucose values. It learns how each user's own meals relate to their glucose, explains its guidance using the user's own history, and keeps trusted people informed through a rule-based emergency escalation chain. This proposal defines the problem, a deliberately narrow graduation scope, the system design, and a twelve-week plan. A proof of concept implements the riskiest core logic in tested, dependency-free Python: a configurable safety rule engine, an escalation state machine, an educational insulin-to-carbohydrate and sensitivity-factor calculator with hard limits, and a personal food-response engine. The engine was evaluated on the public HUPA-UCM dataset of 25 people with type 1 diabetes. With about two weeks of history per patient, a single global average predicted held-out meal responses better than the personal engine (mean absolute error 32.4 versus 36.4~mg/dL, averaged over 22 patients). With months of history the personal engine improved steadily: for one patient its error fell from 29.9 to 20.9~mg/dL, 23\% below the global average. The results suggest that InsuLink should start from population estimates and personalise as a user logs meals. They are a first academic analysis, not clinical evidence.
